\documentclass[11pt, oneside]{book}
%\documentclass[11pt, twoside]{book}  %importante cambiar si es libro

\usepackage[utf8]{inputenc}
\usepackage[T1]{fontenc}


%%%%%%%%%%%%%%%%%%%%%%%%%%%%%%%%%%%%%%%%%%%%%%%%%%%%%%%%%
%   Paquetes y librerías iniciales para un documento,
%           Graficos, simbolos, tablas, bibliografía
%%%%%%%%%%%%%%%%%%%%%%%%%%%%%%%%%%%%%%%%%%%%%%%%%%%%%%%%%
%---------------------------------------------------------------------------
% Esto es para que el LaTeX sepa que el texto está en español:
\usepackage[spanish,es-noshorthands,es-tabla]{babel}
\selectlanguage{spanish}
%---------------------------------------------------------------------------
% para uso en modo matemático
\usepackage{amssymb}
\usepackage{amsthm}
\usepackage{amsmath}

%---------------------------------------------------------------------------
% Para las imagenes

% Sin la opción dvips: con pdflatex (o latexmk -pdf) no hace falta y puede
% dar problemas con las imágenes PDF/PNG
\usepackage{graphicx}
\graphicspath{{Figuras/}} % se fija el camino para las figuras
\usepackage{rotating}  %para poder rotar las figuras
\usepackage{wrapfig}

%---------------------------------------------------------------------------
% PIES DE FOTOS Y TABLAS

\usepackage{float} % para objetos flotantes
\usepackage{caption} % para los pies de fotos o tablas
\usepackage{subcaption}
\captionsetup{font=small,labelfont=bf} % para tamaño y tipo fuentes pie de foto
\usepackage{multirow} % para las tablas unir columnas
\usepackage{multicol} %para escribir en multiples columnas
\usepackage{tabularx} %tablas ajustadas al ancho del área de texto
\usepackage{array}

% Referencias y biblio

\usepackage{pdfpages}
% hyperref se carga al final del preámbulo (justo antes de \begin{document}).
% xurl (que incluye url) se carga aquí porque UAL.sty usa \UrlBreaks (antes
% lo traía hyperref); además deja partir las URL por cualquier carácter, así
% una dirección larga no se sale del margen.
\usepackage{xurl}

% Justificación más fina (pdflatex): microtype ajusta un poco el espaciado y
% el ancho de las letras, y \emergencystretch da un último margen de espacio
% antes de dejar que una línea se salga (rutas en \texttt, palabras largas...)
\usepackage{microtype}
\setlength{\emergencystretch}{3em}

% emptypage solo actúa con twoside (vacía las páginas en blanco que deja
% \cleardoublepage). Con oneside no hace nada: descomentar junto con twoside.
%\usepackage{emptypage}
\usepackage{ragged2e} %Justificacion

%\usepackage[toc,title,page]{appendix}

%Para bibliografía
\usepackage[numbers,square]{natbib}
%\bibliographystyle{apalike}
\bibliographystyle{IEEEtran}

%%%%%%%%% HAbituales estilos
%plainnat
%%abbrvnat
%unsrtnat  %Items in bibliography sorted in order cited)
%IEEEtran

%

%%%%%%%%%%%%%%%%%%%% Formato de la ual %%%%%%%%%%%%%%%%%%%%%%%%%%%%%%%%%


\usepackage{UAL}

% Librerías de TikZ para dibujar los diagramas como vectores (ver la carpeta
% diagramas/ y el ejemplo del capítulo de ejemplos prácticos)
\usetikzlibrary{positioning,arrows.meta,calc,fit,backgrounds,shapes.geometric,shapes.multipart}

% Estilos comunes de los diagramas, en el azul de la plantilla (Color2): así
% todos los diagramas de la memoria tienen el mismo aspecto
\tikzset{
  % --- bloques y capas (arquitectura) ---
  bloque/.style={draw=Color2, fill=Color2!6, rounded corners=2pt, align=center,
                 font=\footnotesize, inner sep=4pt, minimum height=1.0cm},
  subbloque/.style={draw=Color2!60, fill=white, rounded corners=1pt, align=center,
                    font=\scriptsize, inner sep=3pt, minimum height=0.75cm},
  rotulo/.style={font=\footnotesize\bfseries, text=Color2, anchor=north west},
  capa/.style={draw=Color2, fill=Color2!6, rounded corners=2pt, align=left,
               font=\footnotesize, inner sep=5pt},
  flecha/.style={-{Latex[length=2mm]}, draw=Color2, thick},
  flechad/.style={-{Latex[length=2mm]}, draw=Color2, thick, dashed},
  % --- diagramas de secuencia ---
  participante/.style={draw=Color2, fill=Color2!8, rounded corners=1pt, align=center,
    font=\scriptsize, text width=1.78cm, minimum height=0.60cm, inner sep=1pt, line width=0.5pt},
  linea/.style={draw=Color2!45, dashed, line width=0.4pt},
  msg/.style={-{Latex[length=1.7mm]}, draw=Color2, line width=0.55pt},
  ret/.style={-{Latex[length=1.7mm]}, draw=Color2!85, dashed, line width=0.5pt},
  etq/.style={above, fill=white, inner sep=1.2pt, font=\scriptsize, align=center},
  frag/.style={draw=Color2!70, fill=Color2!3, rounded corners=2pt, line width=0.45pt},
  fragtab/.style={draw=Color2!70, fill=white, font=\scriptsize, inner sep=2pt,
    anchor=south west, rounded corners=1pt, line width=0.45pt},
  % --- diagramas de actividad ---
  actividad/.style={draw=Color2, fill=Color2!10, rounded corners=6pt, align=center,
    font=\scriptsize, text width=3.05cm, minimum height=0.55cm, inner sep=3pt, line width=0.5pt},
  decision/.style={draw=Color2, fill=Color2!10, diamond, aspect=2.4, align=center,
    font=\scriptsize, inner sep=1pt, line width=0.5pt},
  flujo/.style={-{Latex[length=1.6mm]}, draw=Color2, line width=0.5pt},
  etqf/.style={font=\scriptsize, fill=white, inner sep=1pt},
  fin/.style={draw=Color2, line width=0.6pt, circle, minimum size=0.22cm, inner sep=0pt},
  % --- estados, entidad-relación, despliegue y componentes ---
  estado/.style={draw=Color2, fill=Color2!8, rounded corners=8pt, minimum width=2.35cm,
    minimum height=0.85cm, align=center, font=\small},
  entidad/.style={draw=Color2, fill=Color2!6, rounded corners=2pt, align=center,
    font=\footnotesize, inner sep=3pt, minimum width=3.35cm},
  rel/.style={draw=Color2!80, line width=0.5pt},
  etqr/.style={font=\scriptsize, fill=white, inner sep=1pt, text=Color2},
  nodo/.style={draw=Color2, fill=Color2!6, rounded corners=2pt, align=center,
    font=\footnotesize, inner sep=4pt, minimum height=1.0cm},
  componente/.style={draw=Color2, fill=Color2!8, rounded corners=2pt, align=center,
    font=\scriptsize, inner sep=3pt, minimum height=0.70cm},
  dep/.style={-{Latex[length=1.5mm]}, draw=Color2!75, line width=0.45pt, dashed},
  % --- diagramas de clases (UML): clase, enumeración y relaciones ---
  clase/.style={rectangle split, rectangle split parts=3, draw=Color2, line width=0.5pt,
    rectangle split part fill={Color2!12,white,white}, font=\scriptsize,
    inner sep=2.2pt, align=left, anchor=north},
  clase/.default={},
  enumeracion/.style={rectangle split, rectangle split parts=2, draw=Color2, line width=0.5pt,
    rectangle split part fill={Color2!12,white}, font=\scriptsize,
    inner sep=2.2pt, align=left, anchor=north},
  enumeracion/.default={},
  asoc/.style={draw=Color2, line width=0.5pt, -{Straight Barb[length=1.6mm,width=2.2mm]}},
  asocbi/.style={draw=Color2, line width=0.5pt, {Straight Barb[length=1.6mm,width=2.2mm]}-{Straight Barb[length=1.6mm,width=2.2mm]}},
  compos/.style={draw=Color2, line width=0.5pt,
    {Diamond[fill=Color2,length=2.8mm,width=1.8mm]}-{Straight Barb[length=1.6mm,width=2.2mm]}},
  gen/.style={draw=Color2, line width=0.5pt, -{Triangle[open,fill=white,length=2.6mm,width=2.8mm]}},
  rol/.style={font=\tiny, text=Color2!90!black, inner sep=1pt},
}

%\usepackage[print,doble]{UAL}
\usepackage{codigo}

%%%%%% PAQUETES INCORPORADOS POR EL USUARIO

\usepackage{longtable}   %Tablas largas
\usepackage{verbatim}
\usepackage{textcomp} %para completar los signos válidos
\usepackage{eurosym} %simbolo del euro [\euro]
\usepackage[bottom]{footmisc}

% Columnas elásticas para que las tablas descriptivas aprovechen el ancho
% disponible entre márgenes y envuelvan el texto sin desbordarse. El número
% es la proporción del ancho de cada columna (deben sumar tantas como
% columnas Y/Z haya): p. ej. {Y{0.6} Y{1.4}} para dos columnas, la segunda
% más ancha. Y alinea a la izquierda y Z a la derecha
\newcolumntype{Y}[1]{>{\hsize=#1\hsize\linewidth=\hsize\RaggedRight\arraybackslash}X}
\newcolumntype{Z}[1]{>{\hsize=#1\hsize\linewidth=\hsize\RaggedLeft\arraybackslash}X}



%%%%%%%%%%%%%%%%%%%%%%%%%%%%%%%%%%%%%%%%%%%%%%%%%%%
% Datos identificativos del TFG                            %
%%%%%%%%%%%%%%%%%%%%%%%%%%%%%%%%%%%%%%%%%%%%%%%%%%%

\author{Nombre Apellido1 Apellido2}
\titulo{Título del Trabajo Fin de Máster}
\subTitulo{DEPARTAMENTO DE
INFORMÁTICA\\  ÁREA DE [ÁREA DE CONOCIMIENTO]}
\estudios{Grado/Máster en Informática}
\universidad{UNIVERSIDAD DE ALMERÍA}
\director{Nombre Apellidos (Director/a)}
\direct{Nombre Apellidos (Director/a 2, si aplica)}
\curso {AAAA/AAAA}


\pagestyle{cab}

\usepackage{enumitem}

% Las listas de figuras y tablas siguen el mismo patrón de las entradas
% de capítulo del índice general: mismo interlineado, parskip y separación
% automática entre capítulos. Se ajusta solo el ancho para las etiquetas.
\makeatletter
\renewcommand*{\l@figure}[2]{%
  \@dottedtocline{1}{1.5em}{2.4em}{#1}{#2}}
\renewcommand*{\l@table}[2]{%
  \@dottedtocline{1}{1.5em}{2.4em}{#1}{#2}}
% El valor por defecto de book para \l@section reserva 2.3em al número, que
% es insuficiente a partir del capítulo 10 ("10.1." ocupa más): el título
% queda pegado al número ("10.1.Visión general"). Se ensancha la caja del
% número y se desplazan en consecuencia los niveles inferiores.
\renewcommand*{\l@section}[2]{%
  \@dottedtocline{1}{1.5em}{3.2em}{#1}{#2}}
\renewcommand*{\l@subsection}[2]{%
  \@dottedtocline{2}{4.7em}{3.6em}{#1}{#2}}
\renewcommand*{\l@subsubsection}[2]{%
  \@dottedtocline{3}{8.3em}{4.3em}{#1}{#2}}
\makeatother

% Los títulos de sección usan \scshape (definido en UAL.sty) y algunos
% contienen \textit. Helvetica no tiene el juego negrita-versalita-cursiva
% (T1/phv/b/scit), así que LaTeX avisaba en cada uno de esos títulos. Se
% declara la sustitución de forma explícita: el resultado visual es el mismo,
% pero el registro de compilación queda limpio. Hay que forzar antes la carga
% del descriptor de la familia (t1phv.fd) o la declaración no encuentra phv.
\AtBeginDocument{%
  \sbox0{\fontfamily{phv}\selectfont x}%
  \DeclareFontShape{T1}{phv}{b}{scit}{<->ssub*phv/b/it}{}%
}

% hyperref debe ser de los últimos paquetes en cargarse, porque redefine
% comandos de otros (footmisc, enumitem, longtable, titlesec...). Antes se
% cargaba dos veces, y la segunda vez anulaba las opciones de la primera.
% Enlaces activos (se puede hacer clic) pero en negro, para imprimir. Para la
% versión digital con enlaces en color, cambiar los tres "black" por "blue".
\usepackage[
breaklinks=true,colorlinks=true,
linkcolor=black,urlcolor=black,citecolor=black,
bookmarks=true,bookmarksopenlevel=2]{hyperref}

\begin{document}
\frontmatter

% La portada se genera automáticamente con los datos definidos arriba
% (\author, \titulo, \director, \direct, \curso) a partir del fondo
% institucional Figuras/logos/TFG_template.pdf
\portada
\begin{dedicatoria}
% La dedicatoria es opcional. Aquí puedes agradecer a quien te haya
% ayudado durante el trabajo: dirección, familia, compañeros, etc.
% Sustituye este texto de ejemplo por el tuyo propio.

Quiero expresar mi agradecimiento a ...

\end{dedicatoria}



%%%%%%%%%%%%%%%%%%%%%%%%%%%%%%%%%%%%%%%%%%%%%%%%%%%%%%%%%%%%%%%%%%%%%%%%%%%%
%%%%%%%%%%%%%%% ESTO ES PARA LA TABLA DE CONTENIDOS (INDICE)%%%%%%%%%%%%%%%%


\setcounter{secnumdepth}{3} %niveles en capitulos hasta (1.1.1)
\setcounter{tocdepth}{3} %niveles en indice hasta (1.1)


\addtocontents{toc}{~\hfill\small{Página}\par} %añade "pagina" a tabla de contanidos
\addtocontents{toc}{\vspace{2pt} \hrule \vspace{5mm} \par}

%\renewcommand\contentsname{Índice de contenidos}
\tableofcontents
\listoffigures % indice de figuras
\listoftables % indice de tablas
\lstlistoflistings % indice de listados: comentarlo si la memoria no incluye
% bloques de código (listings); si no, genera una página vacía titulada
% "LISTADOS"



%---------------------------------------------------------------------------
% comienzo de relacion abreviaturas y/o acrónimos
%---------------------------------------------------------------------------


%\addcontentsline{toc}{section}{ABREVIATURAS}
\clearpage
\vspace{0.2cm}
\section*{ABREVIATURAS}
% Tabla de ejemplo. Sustitúyela por los acrónimos y abreviaturas que
% realmente uses en tu memoria, ordenados alfabéticamente. Cada acrónimo
% debe definirse la primera vez que se usa en el texto, ver capítulo
% "Ejemplos prácticos de la plantilla".
% La tabla ocupa el ancho del texto con columnas elásticas (Y, definidas en
% main.tex) en lugar de \resizebox, que encogía la letra: las descripciones
% largas hacen salto de línea. Se lee por columnas: primero la pareja de la
% izquierda y después la de la derecha.
\noindent
\begin{tabularx}{\textwidth}{|Y{0.4600}|Y{1.5400}|Y{0.4600}|Y{1.5400}|}
\hline
API & \textit{Application Programming Interface} & TFG & Trabajo Fin de Grado \\
ESI & Escuela Superior de Ingeniería & TFM & Trabajo Fin de Máster \\
IDE & \textit{Integrated Development Environment} & UAL & Universidad de Almería \\
IEEE & \textit{Institute of Electrical and Electronics Engineers} & UML & \textit{Unified Modeling Language} \\
PDF & \textit{Portable Document Format} & URL & \textit{Uniform Resource Locator} \\
\hline
\end{tabularx}
%\addcontentsline{toc}{section}{ABREVIATURAS}



% Borra el bloque si sólo queremos resumen en contraportada
\addcontentsline{toc}{chapter}{Resumen y Abstract}

\chapter*{Resumen y Abstract}

\input{resumentexto.tex}


%se suelen utilizar archivos separados, por capitulos para agilizar la compilación en las versiones intermedias
\mainmatter
\marcagua

%%%%%%%%%%%%%%%%%%%%%%%%%%%%%%%%%%%%%%%%%%%%%%%%%%%%%%%%%%%%%%%%%%%%%%%%%%%%
% CONTENIDO DE EJEMPLO DE LA PLANTILLA
% Sustituye este archivo por los capítulos reales de tu TFG/TFM. Se ha
% dejado como guía práctica de los elementos de LaTeX más habituales en
% este tipo de memorias, junto con una guía de estilo y redacción general.
%%%%%%%%%%%%%%%%%%%%%%%%%%%%%%%%%%%%%%%%%%%%%%%%%%%%%%%%%%%%%%%%%%%%%%%%%%%%

\chapter{Ejemplos prácticos de la plantilla}
\label{cap:ejemplos}

Este capítulo reúne, de forma práctica y directamente reutilizable, los elementos de \LaTeX{} que con más frecuencia se necesitan al escribir un TFG o un TFM: acrónimos, figuras, tablas, ecuaciones, listados de código y citas bibliográficas. Puedes copiar el código fuente de \texttt{contenido1.tex} y adaptarlo a tu propio contenido.

\section{Acrónimos y abreviaturas}
\label{sec:acronimos}

La primera vez que se usa un acrónimo en el texto debe escribirse completo, seguido de sus siglas entre paréntesis. Por ejemplo, esta plantilla se edita con un Entorno de Desarrollo Integrado (IDE). En las siguientes apariciones basta con usar la sigla: el IDE empleado permite compilar el documento sin salir del editor.

Además de definir el acrónimo la primera vez que se usa, se recomienda incluir al principio del documento una tabla con todas las abreviaturas empleadas, ordenadas alfabéticamente (véase la sección ``ABREVIATURAS'' al inicio de este documento, definida en \texttt{indicesyrefe.tex}).

\section{Figuras}
\label{sec:figuras}

\subsection{Figura simple}

Para incluir una imagen se usa el entorno \texttt{figure} junto con \lstinline[language=enparrafo]!\includegraphics!. La figura~\ref{fig:arquitectura} muestra un ejemplo de diagrama sencillo.

\begin{figure}[H]
  \centering
  \includegraphics[width=0.75\textwidth]{ejemplo-arquitectura}
  \caption{Ejemplo de diagrama de arquitectura (elaboración propia).}
  \label{fig:arquitectura}
\end{figure}

\subsection{Subfiguras}

Cuando se quieren mostrar varias imágenes relacionadas de forma conjunta se puede usar el paquete \texttt{subcaption}, ya cargado en esta plantilla. La figura~\ref{fig:subfiguras} muestra un ejemplo con dos subfiguras.

\begin{figure}[H]
  \centering
  \begin{subfigure}[b]{0.48\textwidth}
    \centering
    \includegraphics[width=\textwidth]{ejemplo-arquitectura}
    \caption{Arquitectura cliente-servidor.}
    \label{fig:sub-arquitectura}
  \end{subfigure}
  \hfill
  \begin{subfigure}[b]{0.48\textwidth}
    \centering
    \includegraphics[width=\textwidth]{ejemplo-modulos}
    \caption{Módulos del sistema.}
    \label{fig:sub-modulos}
  \end{subfigure}
  \caption{Ejemplo de figura compuesta por dos subfiguras.}
  \label{fig:subfiguras}
\end{figure}

Cada subfigura se referencia igual que una figura normal: la figura~\ref{fig:sub-modulos} muestra los módulos del sistema, mientras que la figura~\ref{fig:subfiguras} agrupa ambas.

\subsection{Diagramas vectoriales con TikZ}

Los diagramas (arquitectura, UML, entidad-relación, secuencia, actividad\ldots) se pueden dibujar directamente en LaTeX con TikZ: se ven nítidos a cualquier tamaño y el texto se puede buscar en el PDF. La plantilla define en \texttt{main.tex} unos estilos comunes en el azul de la portada (\texttt{capa}, \texttt{bloque}, \texttt{flecha}, \texttt{clase}, \texttt{entidad}, \texttt{participante}, \texttt{actividad}\ldots), de modo que todos los diagramas tienen el mismo aspecto. Cada diagrama va en su propio fichero de la carpeta \texttt{diagramas/} y se incluye con \lstinline[language=enparrafo]!\input!, como el de la figura~\ref{fig:tikz-capas}.

Si el pie de una figura es largo, conviene darle una versión corta entre corchetes (\lstinline[language=enparrafo]!\caption[corto]{largo}!): el índice de figuras muestra la corta y la figura, la larga.

\begin{figure}[H]
  \centering
  % Diagrama de ejemplo: arquitectura por capas dibujada en TikZ con los estilos
% comunes de main.tex (capa, flecha). Se incluye con \input desde
% contenido1.tex. Cada diagrama de la memoria va en su propio fichero de esta
% carpeta.

\begin{tikzpicture}[x=1cm,y=1cm]
  \node[capa, minimum width=10cm] (pres) at (0,3.0)
    {\textbf{Presentación}\\ {\scriptsize Aplicación web y móvil}};
  \node[capa, minimum width=10cm] (api) at (0,1.5)
    {\textbf{API / Backend}\\ {\scriptsize Servicios REST y lógica de negocio}};
  \node[capa, minimum width=10cm] (datos) at (0,0)
    {\textbf{Persistencia}\\ {\scriptsize Base de datos}};

  \draw[flecha] (pres) -- (api);
  \draw[flecha] (api) -- (datos);
\end{tikzpicture}

  \caption[Diagrama de capas en TikZ]{Diagrama de capas dibujado en TikZ con los estilos comunes de la plantilla. El índice de figuras muestra solo el título corto.}
  \label{fig:tikz-capas}
\end{figure}

\section{Tablas}
\label{sec:tablas}

\subsection{Tabla simple}

\begin{table}[H]
  \centering
  \begin{tabular}{l l r}
    \hline
    \textbf{Tecnología} & \textbf{Tipo} & \textbf{Año} \\
    \hline
    Ejemplo A & Backend        & 2022 \\
    Ejemplo B & Frontend       & 2023 \\
    Ejemplo C & Base de datos  & 2021 \\
    \hline
  \end{tabular}
  \caption{Ejemplo de tabla simple.}
  \label{tab:simple}
\end{table}

Como toda tabla, se referencia en el texto antes de que aparezca; por ejemplo, la tabla~\ref{tab:simple} resume tres tecnologías de ejemplo.

\subsection{Tabla con celdas combinadas}

El paquete \texttt{multirow}, ya cargado en \texttt{main.tex}, permite combinar celdas verticalmente; \lstinline[language=enparrafo]!\multicolumn! (del núcleo de \LaTeX{}) permite combinarlas horizontalmente. La tabla~\ref{tab:multirow} combina ambos.

\begin{table}[H]
  \centering
  \begin{tabular}{l l c c}
    \hline
    \multirow{2}{*}{\textbf{Módulo}} & \multirow{2}{*}{\textbf{Responsable}} & \multicolumn{2}{c}{\textbf{Estado}} \\
    \cline{3-4}
     & & Diseño & Implementado \\
    \hline
    \multirow{2}{*}{Interfaz} & Persona A & Sí & Sí \\
     & Persona B & Sí & No \\
    Lógica de negocio & Persona C & Sí & Sí \\
    \hline
  \end{tabular}
  \caption{Ejemplo de tabla con celdas combinadas.}
  \label{tab:multirow}
\end{table}

\subsection{Tabla larga (longtable)}

Cuando una tabla puede ocupar más de una página se debe usar el entorno \texttt{longtable} (paquete cargado en \texttt{main.tex}) en lugar de \texttt{table}, ya que este último no puede partirse entre páginas. La tabla~\ref{tab:longtable} es un ejemplo mínimo.

\begin{center}
\begin{longtable}{l l}
  \hline
  \textbf{Elemento} & \textbf{Descripción} \\
  \hline
  \endfirsthead
  \textbf{Elemento} & \textbf{Descripción} \\
  \hline
  \endhead
  Fila 1 & Descripción de ejemplo de la fila 1 \\
  Fila 2 & Descripción de ejemplo de la fila 2 \\
  Fila 3 & Descripción de ejemplo de la fila 3 \\
  \hline
  \caption{Ejemplo de tabla larga con \texttt{longtable}.}
  \label{tab:longtable}
\end{longtable}
\end{center}

\section{Ecuaciones}
\label{sec:ecuaciones}

Las fórmulas sencillas pueden escribirse en línea, por ejemplo $x^2 + y^2 = r^2$. Si conviene numerarlas para referenciarlas después, se usa el entorno \texttt{equation}, como en la ecuación~\eqref{ec:cuadratica}.

\begin{equation}
  \label{ec:cuadratica}
  x_{1,2} = \frac{-b \pm \sqrt{b^2 - 4ac}}{2a}
\end{equation}

\section{Listados de código}
\label{sec:codigo}

Esta plantilla incluye un estilo (\texttt{codigo.sty}) para formatear código fuente con el paquete \texttt{listings}: numeración de línea, resaltado de sintaxis y su propio índice de listados. Los archivos de código de ejemplo se guardan en la carpeta \texttt{Codigo/}.

\lstinputlisting[language=C, caption={Ejemplo de listado en C incluido desde un archivo externo.}, label=lst:holamundo]{holamundo.c}

También se puede escribir el código directamente dentro del documento:

\begin{lstlisting}[language=C, caption={Ejemplo de listado escrito directamente en el documento.}]
#include <stdio.h>

int main() {
    printf("Hola mundo");
    return 0;
}
\end{lstlisting}

El listado~\ref{lst:holamundo} y el anterior producen el mismo resultado; la diferencia está en si el código vive en un archivo aparte (recomendado para código largo o reutilizado) o se escribe en línea. También existe un lenguaje de ejemplo para pseudocódigo, definido en \texttt{codigo.sty}:

\lstinputlisting[language=pseudo, caption={Ejemplo de listado en pseudocódigo.}, label=lst:pseudo]{holamundo.psc}

\section{Citas y bibliografía}
\label{sec:citas}

Las referencias bibliográficas se gestionan con \texttt{natbib} a partir del archivo \texttt{referencias.bib}. Cada entrada del \texttt{.bib} se cita con \lstinline[language=enparrafo]!\cite{clave}! o \lstinline[language=enparrafo]!\citep{clave}!. Por ejemplo, se puede citar un artículo de revista \citep{ejemploarticulo}, un libro \citep{ejemplolibro} o una comunicación en congreso \citep{ejemploconferencia}. El estilo configurado en esta plantilla (\lstinline[language=enparrafo]!\bibliographystyle{IEEEtran}!) es numérico, por lo que las citas aparecen como números entre corchetes; comandos de \texttt{natbib} pensados para estilos autor-año, como \lstinline[language=enparrafo]!\citet!, no deben usarse con este estilo.

Todas las citas insertadas con \texttt{cite}/\texttt{citep}/\texttt{citet} generan automáticamente la lista de referencias que aparece en el capítulo ``BIBLIOGRAFÍA'' al final del documento, ya ordenada y formateada según el estilo indicado en \lstinline[language=enparrafo]!\bibliographystyle{IEEEtran}! (\texttt{main.tex}). El apéndice \ref{sec:bib} resume los tipos de entrada disponibles en un archivo \texttt{.bib} y sus campos.


%%%%%%%%%%%%%%%%%%%%%%%%%%%%%%%%%%%%%%%%%%%%%%%%%%%%%%%%%%%%%%%%%%%%%%%%%%%%
% A partir de aquí: guía de estilo y redacción para TFG/TFM. Es contenido
% genérico y orientativo; consulta siempre las normas concretas de tu
% titulación y las indicaciones de tu director/a.
%%%%%%%%%%%%%%%%%%%%%%%%%%%%%%%%%%%%%%%%%%%%%%%%%%%%%%%%%%%%%%%%%%%%%%%%%%%%

\chapter{Normas fundamentales del formato}
\label{sec:funda}

La memoria del Trabajo Fin de Grado/Máster se presentará en formato  UNE  A4,  salvo  que  las  características  especiales  del  trabajo  presentado  no  lo permitieran.

\section {Estilo,  Tipografía y    Maquetado}

Se podrán utilizar las siguientes familias de tipos de letras \url{https://tug.org/FontCatalogue/}, no se recomiendan tipos de letras con decoración excesiva:

	\begin{center}
	\begin{tabular}{p{2cm}p{7.5cm}c}
		\textbf{Tipo} & \textbf{Fuente} & \textbf{Anchura} \\ \hline
         Familia 1 & Arial, Helvética, Tahoma  	\ldots & Variable \\
         Familia 2 & Times New Roman, Palatino  Linotype,  	\ldots & Variable  \\
         Familia 3 & Courier  New,  Miriam Fixed \ldots & Fija\\
         \hline
	\end{tabular}
	\end{center}



\begin{table}
	\begin{center}
	\begin{tabular}{l c c}
		\textbf{Tipo} & \textbf{Fuente} & \textbf{Anchura} \\ \hline
         Familia 1 & Arial, Helvética, Tahoma  	\ldots & Variable \\
         Familia 2 & Times New Roman, Palatino  Linotype,  	\ldots & Variable  \\
         Familia 3 & Courier  New,  Miriam Fixed \ldots & Fija \\
         \hline
	\end{tabular}
	\end{center}

	\caption{\label{tab:letras}Tipos de letras.}
\end{table}


La familia 1 se recomienda para los títulos de las distintas secciones del documento.

La familia 2 se recomienda para los párrafos normales.

La familia 3 se recomienda para escribir, por ejemplo, porciones de código, debido a su tamaño fijo donde es necesario un texto tabulado y se alinee perfectamente, así como cualquier otro uso de texto de tamaño de letra fijo.

Se recomienda el uso del tipo de letra Palatino Linotype en itálica para la escritura de fórmulas sencillas, para aquellas que sean más complejas se podrá usar un editor de ecuaciones, ejemplos:

$ x^2+y^2=r^2,    a+b=c$

Si fuese necesario se puede utilizar un entorno ecuación para después poder referenciarlo en el texto, como se ha mostrado en la sección~\ref{sec:ecuaciones}. La potencia de \LaTeX{} para la edición de ecuaciones supera el alcance de este documento, se sugiere consultar otras fuentes para ampliar la información sobre su uso.

Las dimensiones del la zona de texto se definen a través del \emph{geometry}, \lstinline[language=enparrafo]!geometry!, un paquete de latex que en este caso se incluye en el estilo \textbf{UAL}, \lstinline[language=enparrafo]!\usepackage{UAL}!.

\section {Estilo ual de latex}

En el estilo se define la portada y contraportada según las recomendaciones de la escuela superior de ingeniería, véase la figura \ref{fig:facultad}.

\begin{figure}
	\begin{center}
		\includegraphics[width = 0.5\textwidth]{facultad.png}
	\end{center}
	\caption{\label{fig:facultad} Logotipo de la Escuela Superior de ingeniería }
\end{figure}



El estilo generado permite definir un documento a una cara o doble cara para lo que se tiene que seleccionar la opción "print" al incluir el paquete estilo (\lstinline[language=enparrafo]!\usepackage[print]{UAL}!). De esta forma se podrá personalizar, el pasar al la página impar en cada sección o bien a la páginas siguiente. Además se fija la marca de agua. Otras cuestiones incorporados en el estilo son los distintos tipos de cabeceras y pies de página necesarios.

Un aspecto importante son los elementos flotantes de latex, en especial tablas y figuras. Estos son los elementos que después pueden incorporarse en los índices de la cabecera del documento. En los casos que se utilicen "códigos" también se puede incluir un índice de listados, \lstinline[language=enparrafo]!\usepackage{listing}!.

Si es necesario se puede trabajar con los paquetes para tablas largas o con aquellos que permiten hacer subfiguras (véase la sección~\ref{sec:figuras} y la sección~\ref{sec:tablas}).

 \section{Figuras y tablas}

Las  figuras  (o  gráficos)  y  tablas  (cuadros)  que  se  incorporen  deben  ir  numerados  (números  arábigos)  y  tener  un  título,  y  siempre  tendrán  que  tener  relación  con el contenido expuesto.  Todo elemento flotante debe referenciarse en el texto y debe colocarse  preferiblemente después de la primera vez que se cita en el texto.

Las  gráficas, fotografías,  los  diagramas,  mapas,  dibujos,  etc.serán  considerados como  figuras. Es importante que se incluya el origen de la figura, en la etiqueta, si la ha hecho el autor se coloca el texto "elaboración propia".

No se deben poner gráficos, figuras y fotos sólo para rellenar siempre se nota demasiado. Los gráficos y figuras deben ir referenciados en el texto, sino es mejor no incorporarlas en la memoria.



 \section{Listados de código}

En el dominio de la informática es habitual que un TFG/TFM deba incluir porciones de código fomateado. Incluso que se definan como un un elemento flotante especial. En \LaTeX{} existen muchas formas de resolver esta cuestión. Se ha definido un estilo llamado \lstinline[language=enparrafo]!codigo.sty! que implementa un ejemplo, véase la sección~\ref{sec:codigo}. Los archivos de código están una carpeta aparte igual que las figuras. \footnote{Se puede sincronizar con repositorios de código tipo \url{www.github.com} se emplea overleaf, pero es necesario mantener dos proyectos}

\subsection{Utilización del entorno verbatim}
Lo que en latex se coloca dentro de un entorno verbatim deja de ser código latex y se representan tal como aparece. Pero suele tener interacciones con el paquete  gráfico tikz. por lo que no es recomendable, es mejor usar \lstinline[language=enparrafo]!\lstinline!

\begin{verbatim}
#include <stdio.h>

int main()
{
        printf("Hola mundo");
        return 0;
}
\end{verbatim}

\subsection{Utilización de paquete listing}


Se puede escribir directamente el código dentro de un entorno y si no se especifica nada se utilizan los parámetros por defecto. Para detalle de como configurar esta utilidad es conveniente revisar la documentación del paquete. Ejemplo de inclusión sin utilizar un objeto flotante. No aparecerá en el índice por poner \lstinline[language=enparrafo]!nolol=true!.

\lstinputlisting[firstline=1, firstnumber=1, nolol=true]{holamundo.c}

Realmente se trata de un paquete muy complejo que se puede complementar con el uso de entornos personalizados, cuya explotación exhaustiva queda fuera de este libro de estilo.


\chapter {Buenas prácticas}


En esta sección se incluyen una serie de cuestiones o buenas prácticas que se derivan de la experiencia, de normas básicas de elaboración de textos científicos \cite{ejemploarticulo}, y de las normas gramaticales básicas.

\begin{enumerate}
    \item Cuidado con el espacio entre palabras, utilizar \emph{hyphenation} para cortar o no palabras con los guiones. Aunque por defecto el paquete de español de latex lo gestiona aceptablemente bien.

\item Los signos de puntuación (comas, puntos, puntos y comas, dos puntos) así   como   el   cierre  de   comillas, paréntesis e   interrogaciones han   de   ir inmediatamente   a   continuación   (sin   espacios)   de   la   palabra   anterior   y separados  por  un  espacio  de  la  palabra siguiente.  De  modo  contrario,  los signos   de   apertura   de   comillas, paréntesis e   interrogaciones deben   ir separados  por  un  espacio  de  la  palabra  anterior e inmediatamente  antes  (sin espacio) de la palabra posterior.


\item En las listas \lstinline[language=enparrafo]!\begin{itemize}! \lstinline[language=enparrafo]!end{itemize}! o enumeraciones \lstinline[language=enparrafo]!\begin{enumerate}! donde no aparezca un verbo, es decir no se a una oración completa no se pone punto al final. Sólo se pone punto en el último item de la lista para finalizar el párrafo.

\item Por la misma razón no se deben acabar los títulos y los subtitulos con un punto. En los pies de figuras y tablas queda a criterio del autor y tutor, pero la recomendación es no ponerlo.

\item Aunque en español las siglas puedan marcarse con puntos entre letra y letra,  se  recomienda  utilizarlas  siempre  en  mayúscula  (siempre  sin  tilde) y  sin los    puntos.

\item La  primera  vez  se emplea una acrónimo debe  ser  definido, incluyendo   el   texto   de   la   definición   entre   paréntesis   y   marcando   con mayúsculas  las  letras  que  dan  lugar  al  acrónimo.  Ejemplo:  MCU  (Multipoint Control  Unit).  En  subsiguientes  apariciones  del  acrónimo  no  debe  volverse  a incluir  la  definición.  Se  recomienda  adjuntar  al  principio  del  proyecto  una relación    detallada    de    los    acrónimos    y    su    significado,    ordenados alfabéticamente.


\item Existen tres tipos de comillas: las españolas o latinas (<<comillas >>), las inglesas ("comillas") y las simples ('comillas' ). Use un único tipo a lo largo del texto. En escritos científicos es más común encontrar las inglesas o las simples. Empléelas siempre que efectúe citas textuales extraídas de algún documento.  También están las comillas abiertas y cerradas ``comillas ''.


\item Evite   emplear   innecesariamente   palabras   o   expresiones   en   inglés cuando  en  español  exista  un  equivalente.  El  uso  de  términos  ingleses  no aceptados debe circunscribirse al campo técnico. Por ejemplo "planning del proyecto" en lugar de "planificación del proyecto".

\item No se deben traducir sin comprobar la validez términos técnicos muy específicos si dichas traducciones no se encuentran más o menos asumidas por la comunidad científica (es decir, las   ha   leído   o   escuchado   repetidas   veces).   En   el   texto   los   términos   y expresiones inglesas no aceptados deben aparecer en cursiva.

\item Cuando  se  referencien  símbolos  matemáticos  en  el  texto,  estos  han  de aparecer   en   cursiva   y   de   la   misma   manera   a   como   aparecen   en   las ecuaciones.

\item No se deben usar frases con el verbo principal en infinitivo y sin articular. Por ejemplo no se debe escribir  “En  primer  lugar,  destacar  que  los  resultados  se  obtuvieron  por simulación”. En su lugar  “En primer  lugar se ha  de  reseñar que  los resultados...” (o “se debe”, “es preciso”, “hemos de”, “cabe” reseñar).

\item El sujeto (salvo que se intercale una frase subordinada) no se puede separar por comas del predicado.

\item Se ha de ser   coherente  con  el  uso  de  los  tiempos  verbales.  Si  describe  lo  que  se  ha hecho en pasado, se ha de continuar  con ese tiempo a lo largo de la memoria.

\item No se debe abusar del  anidamiento  de oraciones  subordinadas.  En  su  lugar  es recomendable  usar  una  secuencia  de  oraciones  cortas,  que  son  de  más  fácil compresión. Asimismo, se deben evitar las negaciones encadenadas
“no  aumentar  el  parámetro  beta  implica  que  no  se  aumenta  la  ganancia”, siendo preferible afirmar “reducir el parámetro beta reduce la ganancia”.

\item Se han de reescribir   las  formas pasivas.  Utilice,  en  su  lugar,  la  pasiva  refleja.  Es preferible “se ha desarrollado un sistema” a “un sistema ha sido desarrollado”. La construcción en pasiva es típicamente inglesa y su abuso denota que estás traduciendo algo (lo cual no es pecado si dices de quien lo traduces). Lo mejor es escribir en castellano. Ejemplo: “Puede  que  sean  retrasadas...”    se  debería  escribir  “Pueden  retrasarse”  o  “puede  que  se retrasen”.

\item Todos los párrafos concluyen con un punto final.

\item En  el  texto,  si se utilizan números se han de expresar   con  palabras   si  son  menores  o  iguales  que veinte y con números si son superiores a ese valor o presentan decimales.

\item A  la  hora  de  emplear  la  coma  o  el  punto  para  expresar  los decimales en los números se debe ser congruente, no usarlos indistintamente.

\item  Los términos  (adjetivos  y  adverbios)  valorativos  que  resulten  exagerados  o poco científicos deben evitarse. Por ejemplo,  “notablemente” o “sustancialmente” antes que “espectacularmente”, “enormemente”, “aplastantemente”, “drásticamente”o “fenomenalmente”. Al hablar de un resultado  o  de  un  producto  de  Ingeniería, es preferible “de forma adecuada o muy conveniente” antes que “perfectamente”.

\item En un texto científico es preferible, para concluir una relación que se sabe que es muy larga, emplear “etc.” que los puntos suspensivos (“...”).

\item El  vocabulario debe ser rico y no repetir términos dentro de la misma frase o  párrafo.  Por  ejemplo,  para comentar  lo  que  una  figura  enseña,  puede  decir “en  la  figura  se  muestra,  se  refleja,  se  indica,  se  ilustra,  se  presenta,  se observa, se detecta, se comprueba, se contempla, se evidencia, se confirma,...”

\item Las  frases han de concordar en  sujeto y predicado.

\item  El término “demostrar” es muy riguroso (salvo  que  exista  una  auténtica demostración). Una figura muestra pero no suele demostrar nada.

\item Salvo en casos muy puntuales (citas textuales, por ejemplo), detrás de los dos puntos (“:”) se sigue en minúsculas.

\item No se debe utilizar la palabra \emph{optimizar} en vano.

\item Escribir  un  trabajo  de  investigación,  no  es  hacer  literatura  barroca.  Cada  frase  y  cada  palabra debe servir para algo. Si no limitamos nuestra verborrea, nos podemos meter en un lío.

\item Si  se  va  a  hacer  una  descripción  de,  por  ejemplo,  siete  etapas  de  algo,  no  está  de  más, primero listarlas y luego iniciar la descripción.



\end{enumerate}

 Otra lista de errores y/o consideraciones se muestra en la figura \ref{fig:errores}.

\begin{figure}[H]
	\begin{center}
		\includegraphics[width = 0.95\textwidth]{Figuras/errores}
	\end{center}
	\caption{\label{fig:errores} Errores habituales}
\end{figure}




 \chapter{Apartado BIBLIOGRAFIA}
 Este apartado es de vital importancia para tener éxito en la consecución del trabajo fin de grado. Los tribunales suelen prestar gran atención a este apartado, probablemente por su formación académica.

 Es todo un mundo el de los estilos bibliográficos y la realización de las citas, se puede encontrar un tutorial en:  \url{http://ci2.ual.es/comunicar-la-informacion/citas-y-referencias-bibliograficas/}. No obstante en las siguientes secciones se indicarán las cuestiones básicas para iniciarse en el tema pero restringuiendose a la utilización de la presente plantilla. La gran ventaja de latex es que cada elemento en las referencias   \lstinline[language=enparrafo]!\bibitem! es considerado como un elemento de una base de datos, de forma que es el compilador el que reordena o coloca la cita de acuerdo con el estilo definido.

 \section{Como citar}
A lo largo del texto de la memoria se incluirán citas a trabajos de otros autores. La inclusión de la cita tiene que ser justo después del texto donde se propone algo relativo al trabajo citado. Para citar, como para todo el \LaTeX{} existen numerosos paquetes, los dos mas conocidos son  biblatex y   natbib .

Además dada forma de citar y de incorporar items en el apartado referencias lleva asociado un estilo de bibliografía    \lstinline[language=enparrafo]!\bibliographystyle{XXX}!, que aun abre mucho mas el abanico de opciones.

Es recomendable consultar las ayudas de los paquetes \url{https://www.overleaf.com/learn/latex/Biblatex_bibliography_styles} y \url{https://www.overleaf.com/learn/latex/Bibliography_management_with_natbib}.

Lo habitual es utilizar la orden      \lstinline[language=enparrafo]!\cite{xxx}!, o   \lstinline[language=enparrafo]!\citep(xxx)!  en la ubicación en el texto donde se quiere colocar la cita, el compilador se encargá de procesarla según el estilo, tal como se ha mostrado en la sección~\ref{sec:citas}.

 \section {Apartado referencias}
  Existen dos formas de incorporar las referencias, directamente en el entorno   \lstinline[language=enparrafo]!\begin{bibliography}! y utilizando las ordenes del lenguaje o bien utilizando un archivo .bib para guardar esas referencias. Se generará un auxiliar .bbl que automáticamente se incorpora en el documento.

  Este archivo auxiliar es el que tiene directamente las ordenes del lenguaje. Un ejemplo simplificado de su estructura sería:

 \begin{verbatim}

\begin{thebibliography}{3}

\bibitem{ejemplolibro}
N.~Apellido, \emph{Título de ejemplo de un libro}, 2nd~ed.\hskip 1em plus
  0.5em minus 0.4em\relax Editorial de Ejemplo, 2022.

\bibitem{ejemploinforme}
{Organización de Ejemplo}, ``Título de ejemplo de un informe técnico,''
  Institución de Ejemplo, Tech. Rep., 2020.

\bibitem{ejemploarticulo}
N.~Apellido1 and N.~Apellido2, ``Título de ejemplo de un artículo de
  revista,'' \emph{Nombre de la Revista}, vol.~10, no.~2, pp. 123--134, 2023.
\end{thebibliography}

  \end{verbatim}



 \section{Archivos .bib}

 BibTeX, y natbib por extensión, usan un archivo externo en texto plano como base de datos de referencias bibliográficas para generar las bibliografías y sus referencias en documentos con distintos formatos de artículos, libros, tesis, presentaciones, etc. Los nombres de archivos de referencias bibliográficas de BibTeX usualmente terminan o usan la extensión .bib. Los ítems bibliográficos incluidos en un .bib están separados por tipos.

 \begin{verbatim}
@inbook{Berander2005,
        address = {Berlin, Heidelberg},
        author = {Berander, Patrik and Andrews, Anneliese},
        booktitle = {Engineering and Managing Software Requirements},
        pages = {69--94},
        publisher = {Springer Berlin Heidelberg},
        title = {{Requirements Prioritization}},
        doi= {https://doi.org/10.1007/3-540-28244-0{\_}4},
        year = {2005}
}
 \end{verbatim}

 Los tipos que son reconocidos por virtualmente todos los estilos de BibTeX se muestran en el anexo \ref{sec:bib}. Pero igual que antes es recomendable recurrir a tutoriales o recursos donde se describan con detalle. Además se recomienda la utilización de algún gestor de referencias bibliográficas tipo Mendeley. \url{https://www.mendeley.com/}


% La bibliografía como capítulo del índice. El \cleardoublepage +
% \phantomsection son necesarios: \bibliography abre página nueva, así que
% sin ellos la entrada del índice apuntaba a la última página del capítulo
% anterior y quedaba anidada como sección de ese capítulo.
 \cleardoublepage
 \phantomsection
 \addcontentsline{toc}{chapter}{Bibliografía}
 \bibliography{referencias}

\appendix
%\clearpage

\input{apendices}


 \clearpage
 \backmatter
% La contraportada se genera automáticamente a partir del fondo
% institucional Figuras/logos/TFG_back.pdf y del resumen (resumentexto.tex)

\thispagestyle{empty}
\begin{tikzpicture}[remember picture,overlay]
\node (f) [rectangle]  
at (current page.center)
          {\parbox[b][\paperheight]{\paperwidth}{
           \includegraphics[width=\paperwidth,height=\paperheight,keepaspectratio]{Figuras/logos/TFG_back}
          }};
\node (texto) [rectangle,anchor=north west,text width=0.70\paperwidth]
at ([xshift=5.5cm,yshift=-2.2cm]f.north west)
          {\parbox{0.70\paperwidth}{\justifying \small \color{white}
         \input{resumentexto}
         }};
       \node[anchor=base west] at (12.87,-24.63) {\Large\curso};
  \end{tikzpicture}

\end{document}
